\documentclass[11pt]{article}
\usepackage[OT1]{fontenc}
\renewcommand{\rmdefault}{cmr}
\renewcommand{\sfdefault}{cmss}
\renewcommand{\ttdefault}{cmtt}
\usepackage{amsmath,amssymb,amsthm,mathtools}
\usepackage[margin=27mm]{geometry}
\usepackage{microtype,booktabs,longtable,array}
\usepackage[hidelinks]{hyperref}
\setlength{\emergencystretch}{2em}
\setlength{\parskip}{3pt plus 1pt}
\setcounter{tocdepth}{1}
\newtheorem{theorem}{Theorem}[section]
\newtheorem{lemma}[theorem]{Lemma}
\newtheorem{proposition}[theorem]{Proposition}
\newtheorem{corollary}[theorem]{Corollary}
\theoremstyle{definition}
\newtheorem{definition}[theorem]{Definition}
\newtheorem{example}[theorem]{Example}
\theoremstyle{remark}
\newtheorem{remark}[theorem]{Remark}
\newcommand{\N}{\mathbb N}
\newcommand{\Q}{\mathbb Q}
\newcommand{\R}{\mathbb R}
\newcommand{\Z}{\mathbb Z}
\newcommand{\supp}{\operatorname{supp}}
\newcommand{\Con}{\operatorname{Con}}
\newcommand{\ord}{\operatorname{ord}}
\newcommand{\dom}{\operatorname{dom}}
\newcommand{\sgn}{\operatorname{sgn}}
\newcommand{\res}{\operatorname{Res}}
\newcommand{\rad}{\operatorname{rad}}
\newcommand{\CK}{\mathsf{CK}}
\newcommand{\KT}{\mathsf{KT}}
\newcommand{\DCB}{\mathsf{DCB}}
\newcommand{\FS}{\mathsf{FS}}
\newcommand{\Ag}{\mathsf{Ag}}
\newcommand{\Know}{\mathsf K}
\newcommand{\Truth}{\mathsf T}
\newcommand{\Prv}{\mathsf{Pr}}
\newcommand{\code}[1]{\ulcorner #1\urcorner}
\newcommand{\Tot}{\operatorname{Tot}}
\numberwithin{equation}{section}

\title{The Exact Computability of Two\\Adversarial Prediction Measures}
\author{}
\date{28 October 2026\\Working manuscript}
\begin{document}
\maketitle
\begin{abstract}
We determine the oracle complexity of Hibbard's original intelligence
measure in the representation used by Alexander and Hibbard: total
binary predictors are given by ordinary Turing-machine indices, and
evaders are grouped by eventual strict worst-case runtime domination.
The extended-natural-valued measure and its strict comparator have
exact Turing degree $0'''$. Every sufficiently high fixed threshold,
and the predicate of infinite measure, is $\Pi^0_3$-complete on the
total-predictor promise. We then replace the primitive-recursive list
by the canonical repetition-free list of all total computable functions,
ordered by least program index. Every resulting measure value is finite,
but the measure still has exact degree $0'''$. Thus this specified
all-total-function variant does not require a stronger oracle than
the original measure. A robust prediction lemma handles the feedback
created by modified predictions; a uniform threshold formula prevents
an unnecessary extra jump when the list itself is noncomputable.
\end{abstract}
\tableofcontents

\section{Intelligence measures as represented mathematical functions}

In adversarial sequence prediction, a predictor tries to match the
next bit chosen by an evader. Each side sees the other side's previous
outputs. Hibbard's measure assigns a predictor a level according to
the runtime classes of evaders it eventually learns. Alexander and
Hibbard \cite{AH} separate this underlying idea from the particular
enumeration used to organize growth rates. After proving
noncomputability, they ask for more precise degrees and conjecture
that replacing primitive recursive functions by all total computable
functions would require a strictly stronger oracle.

These questions depend on representation. A predictor is a total
function, but an input to an algorithm is a program index. A list
of all total computable functions cannot have a computable total
evaluator, although each function in it is individually computable.
The measure can also take the extra value $\infty$. We make each of
these choices explicit.

Two main results follow. The original measure requires exactly the
third Turing jump. A specified canonical all-total-function variant
requires the same jump, even though its enumeration is noncomputable
and it has no infinite values. The latter answers a precise instance
of the proposed strict increase negatively; it does not classify
every possible ordering of the total computable functions.

\section{The game, the machines and the measures}

Fix an acceptable numbering of ordinary Turing programs on finite
binary strings, with an effective bounded simulation. A total binary
program computes a map $\{0,1\}^{<\omega}\to\{0,1\}$. We use the
standard input-tape convention in which a program can inspect its
first $k$ input cells and stop without scanning the remaining input.
There is no mandatory input-reading cost proportional to its length.

Given a total predictor $p$ and total evader $e$, define their unique
play by
\[
 y_n=p(x_0\ldots x_{n-1}),\qquad
 x_n=e(y_0\ldots y_{n-1}).
\]
Both empty histories occur at round zero. The evader sees the
\emph{predictions}, not its own previous evasions. The predictor
\emph{learns} $e$ if $x_n=y_n$ for every sufficiently large $n$.
For a total evader, let $t_e(n)$ be its maximum running time on all
$2^n$ binary inputs of length $n$.

For an ordinary total computable bound $f:\N\to\N$, put
\[
 E_f=\{e:e\text{ is total binary and }t_e(n)<f(n)
                    \text{ for all sufficiently large }n\},
\]
and write $T(p,f)$ for the assertion that $p$ learns every $e\in E_f$.
The strict inequality is part of the contract. A finite initial segment
of exceptions is allowed.

Let $g_1,g_2,\ldots$ be a listing of growth functions and define
\begin{equation}\label{eq:bounds}
 f_m(n)=\max_{1\le i\le m,\ 0\le j\le n}g_i(j),\qquad m\ge1.
\end{equation}
The bounds are pointwise nondecreasing in $m$. Define
\[
 H_g(p)=
 \begin{cases}
 0,&\text{if }T(p,f_1)\text{ fails},\\
 \max\{m:T(p,f_m)\},&\text{if this set is nonempty and bounded},\\
 \infty,&\text{if every }T(p,f_m)\text{ holds}.
 \end{cases}
\]
The initial-segment property makes these cases exhaustive.

For the original measure, use the fixed effective primitive-recursive
listing of \cite[Definition~7]{AH}. What is needed here is an effective
total evaluator for that listing, its coverage of all primitive
recursive functions, and hence coverage of $n\mapsto n+1$.
Call its measure $H_{\mathrm{PR}}$.

All oracle computations of a measure below have the promise that
their input index is a total binary predictor. They must halt and
return the right answer on that domain. A machine computing an
extended-natural-valued measure outputs a natural number or a distinct
symbol for $\infty$. Statements of exact degree use this representation.

\subsection{The complexity statements}

Write $0',0'',0'''$ for successive Turing jumps. We use the standard
arithmetical-hierarchy fact that a recursive-matrix $\Pi^0_3$ predicate
is decidable in $0'''$, and that complete such predicates compute
$0'''$.

\begin{theorem}\label{thm:original}
There is a fixed coding-dependent $m_0$ such that:
\begin{enumerate}
\item for every fixed $m\ge m_0$, the predicate
      $H_{\mathrm{PR}}(p)\ge m$ is $\Pi^0_3$-complete on total
      predictor indices;
\item $H_{\mathrm{PR}}(p)=\infty$ is $\Pi^0_3$-complete on the same
      promise;
\item $H_{\mathrm{PR}}$ and the strict comparison relation between
      two such values have exact oracle degree $0'''$.
\end{enumerate}
The reductions witnessing hardness uniformly produce total predictors.
\end{theorem}

\begin{theorem}\label{thm:total}
List all total computable functions without repetitions in order of
their least program indices, and use \eqref{eq:bounds}. For this
canonical list, every $H_g(p)$ is finite and $H_g$ has exact oracle
degree $0'''$ on total predictor indices. Every sufficiently high
fixed threshold is $\Pi^0_3$-complete on that promise.
\end{theorem}

The hardness statements for fixed thresholds in the second theorem
are nonuniform in the list position: a fixed position has a fixed
ordinary bound program which can be hardcoded. The theorem does not
assert that these bound indices can be found computably from their
positions.

\section{A uniform \texorpdfstring{$\Pi^0_3$}{Pi-0-3} threshold formula}

\begin{lemma}\label{lem:membership}
Uniformly in an ordinary program index for a total computable $f$,
membership $e\in E_f$ has a recursive-matrix $\Sigma^0_2$ definition.
\end{lemma}
\begin{proof}
An existential witness consists of a cutoff $N$ and a common finite
halting certificate for all inputs of length at most $N$. Check that
these computations halt with binary outputs. For each longer input
$x$ and each stage $s$, require: if the $f$ program on $|x|$ has
halted by $s$ with value $v$, then the $e$ program on $x$ halts with
a binary output in fewer than $v$ steps. This is a bounded decidable
test. Since $f$ is total, the guard eventually activates at each input
length. The universal tests therefore impose exactly eventual strict
runtime domination and total binary output beyond the cutoff; the
finite certificate covers all remaining inputs. Thus
\[
 e\in E_f\quad\Longleftrightarrow\quad
 \exists a\,\forall b\ R(e,f,a,b)
\]
for a recursive predicate $R$, on the total-$f$ promise.
\end{proof}

\begin{lemma}\label{lem:learning}
There is a recursive predicate $S$ such that, for total binary $p,e$,
$p$ learns $e$ exactly when
$\exists c\,\forall d\ S(p,e,c,d)$.
\end{lemma}
\begin{proof}
Let $c$ be an error cutoff. The number $d$ may encode a finite play,
halting times for all its computations, and a designated error round
at least $c$. Reject exactly when this is a valid computation
certificate for such an error. Validity is decidable by bounded
simulation. Every actual error has a finite certificate, and every
valid certificate describes an actual initial play by determinism.
Totality ensures the actual play is infinite. Absence of all
post-cutoff error certificates is therefore eventual learning.
\end{proof}

\begin{proposition}\label{prop:threshold}
Uniformly in total binary $p$ and an ordinary total bound program $f$,
$T(p,f)$ has a recursive-matrix $\Pi^0_3$ definition.
\end{proposition}
\begin{proof}
Using the two preceding lemmas, write
\begin{equation}\label{eq:pi3}
 \forall e,a\ \exists b,c\ \forall d\
 \bigl(\neg R(e,f,a,b)\ \lor\ S(p,e,c,d)\bigr).
\end{equation}
If $a$ does not certify membership in $E_f$, one counterexample $b$
makes the disjunction true. If it does, a learning witness $c$ is
required. Partial or nonbinary evaders cannot satisfy the membership
antecedent, so their behavior in the learning formula is irrelevant.
This proves equivalence and the claimed uniformity.
\end{proof}

The uniformity matters later. Computing an ordinary index for $f$
with an oracle and then applying \eqref{eq:pi3} does not require
three jumps \emph{relative} to that oracle. The formula already has
a recursive matrix with the ordinary index as a parameter.

For the primitive-recursive listing, the $f_m$ indices are computable
uniformly in $m$. Thus all thresholds are uniformly $\Pi^0_3$.
Moreover
\[
 H_{\mathrm{PR}}(p)=\infty\quad\Longleftrightarrow\quad
 \forall m\ T(p,f_m)
\]
is still $\Pi^0_3$: the extra universal quantifier joins the first
block. Oracle $0'''$ first decides this predicate. In the finite case
it tests thresholds until the first failure and returns one less
than that position. This proves the upper bound for the original
measure and, by computing two values, for its comparator.

\section{Prediction along the actual history}

The lower bound will deliberately flip a good baseline's predictions.
This changes the evader's future input. It is therefore insufficient
to compare the new predictor with a good predictor's unmodified play.
We need a bounded prediction rule whose correctness is independent
of how the past predictions were chosen.

Fix a total computable time budget $b$. Given actual paired histories
$x_{<n}$ and $y_{<n}$, examine program indices $i\le n$. A candidate
is eligible when:
\begin{enumerate}
\item at each earlier round $r<n$, if its computation on $y_{<r}$
      halts in fewer than $b(r)$ steps, its output equals $x_r$;
\item on $y_{<n}$ it halts in fewer than $b(n)$ steps with a binary
      output.
\end{enumerate}
Choose the least eligible index and predict its output, or predict
zero if none is eligible. Call this rule $B_b$. A timeout in an
earlier test is ignored; it does not disqualify a candidate.

\begin{lemma}[Robust baseline]\label{lem:robust}
If $e$ is a total binary evader with $t_e(n)<b(n)$ eventually, then
$B_b$ predicts $e(y_{<n})$ correctly at all but finitely many rounds,
regardless of the sequence of actual past predictions $y$.
\end{lemma}
\begin{proof}
Let $k$ be an index of $e$. Once $n$ is large enough, $k\le n$ and
the current computation of $e$ is within budget. Every earlier
bounded computation of $e$ agrees with the recorded $x_r$, since
$x_r=e(y_{<r})$ in the actual play. Thus $e$ is eligible from then on,
and the least eligible index is at most $k$.

If a selected index $i<k$ gives a wrong output at round $n$, that
bounded computation on the recorded input $y_{<n}$ permanently
contradicts the earlier-round test at every later round. It can
never be selected again. There are finitely many indices below $k$.
After their possible first wrong selections, every selected output
is correct. The proof never assumed that the actor actually used
the baseline bit as its own prediction.
\end{proof}

To implement this rule in an ordinary predictor whose input is just
$x_{<n}$, compute its own previous outputs on the strictly shorter
prefixes $x_{<r}$, $r<n$. They reconstruct $y_{<n}$. Structural
recursion on input length terminates. This may be computationally
expensive, but it defines a total program, which is all the degree
argument requires. It remains total if a decidable rule flips its
current baseline output.

\subsection{An infinite-level predictor for the original list}

For the primitive-recursive hierarchy set
\[
 b_\infty(n)=f_{n+1}(n).
\]
It is an ordinary total computable function, since the list has an
effective total evaluator. It need not itself be primitive recursive.
For every fixed $m$, $b_\infty(n)\ge f_m(n)$ eventually. The predictor
which recursively reconstructs its own history and always uses
$B_{b_\infty}$ therefore learns every evader in every $E_{f_m}$.
Denote this fixed total predictor by $q_\infty$; it satisfies
\begin{equation}\label{eq:qinf}
 H_{\mathrm{PR}}(q_\infty)=\infty.
\end{equation}

\section{Coding \texorpdfstring{$\Pi^0_3$}{Pi-0-3} information into finite and infinite errors}

Choose a fixed list position $m_0$ containing the function $n\mapsto n+1$.
Then $f_m(n)\ge n+1$ for every $m\ge m_0$. For each $k\in\N$, let
$e_k$ ignore input bits and output 1 if input length is less than $k$,
and 0 otherwise. It reads at most $k$ cells, so its runtime is bounded
by a constant $c_k$, depending on $k$. Consequently
\begin{equation}\label{eq:codeevader}
 e_k\in E_{f_m}\quad\text{for every }k\text{ and }m\ge m_0.
\end{equation}
Its evasion stream is $1^k0^\omega$, whatever predictions it receives.
No runtime bound uniform in $k$ is claimed.

Let $B\subseteq\N$ be any $\Pi^0_3$ set. Use a recursive normal form
\[
 z\in B\quad\Longleftrightarrow\quad
 \forall k\ \exists s\ \forall t\ R_0(z,k,s,t).
\]
Uniformly produce computably enumerable sets $W_{z,k}$ such that
\begin{equation}\label{eq:finite}
 z\in B\quad\Longleftrightarrow\quad
 \forall k\ (W_{z,k}\text{ is finite}).
\end{equation}
Here is a direct construction. At stage $t$, take the least $s\le t$
passing all tests $R_0(z,k,s,u)$ for $u\le t$, or use $t+1$ if none
does. Enumerate a fresh number whenever this candidate changes. If
there is a least true witness, all smaller candidates eventually fail
and the least true witness stabilizes. If no true witness exists,
the candidate eventually exceeds each fixed bound and changes
infinitely often. This proves \eqref{eq:finite}.

\begin{proposition}\label{prop:hardness}
For any fixed ordinary total bound $b$ eventually exceeding each
$t_{e_k}$, the predicate $T(p,b)$ is $\Pi^0_3$-hard on total
predictors. If $b=b_\infty$, the same reduction makes membership in
$B$ equivalent to infinite original measure, and equivalent to every
fixed threshold $m\ge m_0$.
\end{proposition}
\begin{proof}
Build $p_z$ as the robust $B_b$ rule using its own reconstructed
prediction history, except for this flip: if $x_{<n}$ begins with a
completed code $1^k0$, flip the baseline bit precisely when the
stage-$n$ enumeration of $W_{z,k}$ adds a new element. The first zero
determines at most one $k$. The code is permanent once completed.
The program is total by structural recursion and is produced
computably from $z$, with the fixed ordinary bound program hardcoded.

If $z\in B$, every selected set $W_{z,k}$ is finite. On any play
there are therefore finitely many flips; a stream of all ones never
completes a code and has none. Against any $e\in E_b$, the robust
baseline is eventually correct even with these modified predictions.
Thus $T(p_z,b)$ holds.

If $z\notin B$, choose $k$ with $W_{z,k}$ infinite. Against $e_k$,
the stream completes precisely the code $1^k0$. The baseline is
eventually correct by the budget assumption, but there are infinitely
many subsequent flips. They give infinitely many genuine errors,
so $T(p_z,b)$ fails.

For $b=b_\infty$, a positive input is learned at every original
level by the same robust argument, giving infinite measure. A
negative input fails against $e_k\in E_{f_{m_0}}$, so its measure is
less than $m_0$. This proves all the further equivalences.
\end{proof}

The threshold upper bounds and Proposition~\ref{prop:hardness} prove
the first two parts of Theorem~\ref{thm:original}. Reading a fixed
complete threshold from the value proves the lower bound
$0'''\le_T H_{\mathrm{PR}}$.
For the strict comparator, combine \eqref{eq:qinf} with the reduction:
\[
 H_{\mathrm{PR}}(p_z)<H_{\mathrm{PR}}(q_\infty)
 \quad\Longleftrightarrow\quad z\notin B.
\]
An oracle for this relation decides a complete $\Pi^0_3$ set by
complementing its answer. Together with the upper bounds this proves
the remaining assertions. It does not require a predictor of every
exact finite measure value.

\section{The canonical list of all total computable functions}

We now change the list, keeping the game, runtime convention, and
formula for $H_g$ unchanged. Fix an acceptable numbering $\varphi_e$
of partial functions $\N\rightharpoonup\N$. Scan indices in increasing
order, retain an index when it is total and its function differs
from every earlier retained function. Let $a_m$ be the $m$th
retained index and $g_m=\varphi_{a_m}$.

\begin{lemma}\label{lem:list}
The sequence $m\mapsto a_m$ is computable in $0''$. Each entry is
an ordinary total computable function, and every such function
occurs exactly once. Oracle $0''$ uniformly supplies an ordinary
program index for $f_m$.
\end{lemma}
\begin{proof}
Totality is a $\Pi^0_2$ property decided by $0''$. For two promised
total programs, search for an argument where their outputs differ.
This ordinary search halts exactly when they differ, so equality
is decidable with $0'$. Each finite comparison with previously
retained entries is therefore available in $0''$.
The scan keeps exactly the least index of each total function.
There are infinitely many distinct total functions, so each
requested position is reached. Given the finitely many ordinary
indices $a_1,\ldots,a_m$, compile an ordinary program taking the
finite maximum in \eqref{eq:bounds}. It is total. Only the choice
of its index uses the oracle; its evaluation does not.
\end{proof}

\begin{lemma}\label{lem:finitevalues}
For any listing of all total computable functions, every total
predictor has a finite $H_g$ value.
\end{lemma}
\begin{proof}
Given total $p$, recursively define $z_n=1-p(z_{<n})$.
Let $e_p$ ignore its input bits and output $z_n$ on inputs of length
$n$. This is a total binary evader, and along the actual play $p$
is wrong at every round. Its worst-case runtime $t_{e_p}(n)$ is
total computable: run it on the finitely many length-$n$ inputs
and take the maximum halting time. The function $t_{e_p}+1$
therefore appears as some $g_j$. For every $n$,
$f_j(n)\ge t_{e_p}(n)+1$, so $e_p\in E_{f_j}$. This threshold
fails, and monotonicity makes every later threshold fail too.
\end{proof}

\begin{proof}[Proof of Theorem~\ref{thm:total}]
To compute $H_g(p)$, use $0''\le_T0'''$ to obtain each ordinary
$f_m$ index from Lemma~\ref{lem:list}, and use the uniform formula
of Proposition~\ref{prop:threshold} with oracle $0'''$ to decide
$T(p,f_m)$. Search for the first failure. Lemma~\ref{lem:finitevalues}
guarantees termination and gives the answer $m-1$. This proves
$H_g\le_T0'''$.

Choose the fixed list position $m_0$ of $n\mapsto n+1$.
For every fixed $m\ge m_0$, all code evaders $e_k$ belong to
$E_{f_m}$, by the same constant-runtime argument as
\eqref{eq:codeevader}. Hardcode the ordinary program for this
fixed $f_m$ into Proposition~\ref{prop:hardness}. It reduces an
arbitrary $\Pi^0_3$ set to $T(p,f_m)$. The upper formula remains
$\Pi^0_3$ for this fixed program, so the threshold is complete.
Reading one such threshold from $H_g$ proves the lower degree
bound. All these reductions output ordinary total predictors.
\end{proof}

The noncomputable enumeration therefore does not force an extra
jump in the measure. Its oracle is used to select an ordinary
bound program; the threshold problem for that program was already
uniformly third-level. Confusing this composition with a third
jump \emph{relative to} $0''$ would give an unnecessarily large
upper bound.

More generally, if a listing supplies ordinary total function
indices computably in an oracle $E$, the same argument gives
\[
 H_g\le_T E\oplus0'''.
\]
This is an upper bound, not a claim that $H_g$ can recover $E$.
The finite-value lemma still applies whenever the list covers
all total computable functions.

\section{What is classified, and what remains representation dependent}

The two measures have equal oracle degree but different ranges.
The original list admits the infinite-level predictor $q_\infty$;
every value for the all-total-function list is finite. Degree alone
does not erase this structural distinction.

The original source conjecture attributes a possible oracle increase
to the noneffective total-function list. The specified canonical
listing above shows that this increase is not automatic: the
measure remains third-jump computable and third-jump hard.
An arbitrary noneffective reordering is not covered by its upper
bound. Nor does a lower bound for a finite-valued measure by itself
give the same lower bound for its two-input strict comparator;
the original-measure comparator proof used an infinite-valued
calibration predictor unavailable in the new list.

The runtime convention is also substantive. The code evaders use
a constant amount of work for each fixed $k$. A model imposing a
mandatory full-input scan would need a different common bound
before the hardness proof could be transferred. Totality promises
and the ordinary program-index representation should likewise remain
visible when quoting the result.

Finally these are computability results for specified adversarial
prediction measures. They do not claim that third-jump calculations
can be implemented as practical intelligence tests, or settle the
ordinal, hyperreal and other taxonomies considered in \cite{AH}.

\paragraph{Preparation.}
This AI-assisted hand-proof manuscript includes the adaptive-history
lemma, quantifier calculations and lower-bound constructions. Standard
arithmetical-hierarchy and jump facts are background inputs. No
formal certification or established priority is claimed.

\begin{thebibliography}{9}
\bibitem{AH} S.~Alexander and B.~Hibbard,
\emph{Measuring Intelligence and Growth Rate: Variations on Hibbard's
Intelligence Measure},
\href{https://arxiv.org/abs/2101.12047v1}{arXiv:2101.12047v1},
25 January 2021. Especially Definitions~2--7, Propositions~8--9 and
the complexity discussion following Proposition~9.
\end{thebibliography}
\end{document}
